% GENERATED FILE. Edit canonical lessons, then run npm run generate:books.
% canonical-source-hash: fnv1a64:8e1446b8d833b86b
% canonical-lessons: KA-C402-bhrastacara, KA-C402-lanca, KA-C402-paradarsaka, KA-C402-maremacu, KA-C402-sipharasu-madu

\chapter{Bribes and Secrets}
\label{ch:ka-bribes-and-secrets}
\hlchaptermodality{402}

\begin{chapteropening}
\textbf{By the end of this chapter:} \emph{I can say corruption, a bribe, transparent, to conceal and to recommend.}

Complete the last lesson of chapter 402: I can say corruption, a bribe, transparent, to conceal and to recommend.
\end{chapteropening}

\section[\texorpdfstring{\kn{ಭ್ರಷ್ಟಾಚಾರ} (bhraṣṭācāra)}{ಭ್ರಷ್ಟಾಚಾರ (bhraṣṭācāra)}]{\kn{ಭ್ರಷ್ಟಾಚಾರ} (bhraṣṭācāra) — corruption}

\label{lesson:KA-C402-bhrastacara}



\begin{quote}
Before the new one: say the Kannada for to promote, then the Kannada for to invest.
\end{quote}

\subsection*{You'll want to know: \kn{ಭ್ರಷ್ಟಾಚಾರ}}
\textbf{\kn{ಭ್ರಷ್ಟಾಚಾರ}} — \emph{bhraṣṭācāra} — \textquotedblleft{}corruption\textquotedblright{}.

\begin{quote}
\kn{ಭ್ರಷ್ಟಾಚಾರದ} \kn{ವಿರುದ್ಧ} \kn{ಜನರು} \kn{ಬೀದಿಗೆ} \kn{ಇಳಿದರು}.
\end{quote}

\begin{quote}
\emph{bhraṣṭācārada viruddha janaru bīdige iḷidaru.}
\end{quote}

\begin{quote}
The people took to the streets against corruption.
\end{quote}

\kn{ಭ್ರಷ್ಟ} (fallen) + \kn{ಆಚಾರ} (conduct).

\begin{cousinweb}
\textbf{\kn{ಭ್ರಷ್ಟಾಚಾರ}} (\emph{bhraṣṭācāra}) is built from Sanskrit \emph{bhraṣṭa}, \textquotedblleft{}fallen\textquotedblright{} (CDIAL 9655), and \emph{ācāra}, \textquotedblleft{}conduct\textquotedblright{} (CDIAL 1071): fallen conduct.
\end{cousinweb}

\begin{grammarlens}[title={What you've built}]
One more word for giving reasons and opinions.
\end{grammarlens}

\subsection*{Your turn}
\begin{itemize}
\raggedright
  \item \emph{Say it:} \emph{bhraṣṭācāra}
  \item \emph{Say it:} \emph{bhraṣṭācāra}, once more
  \item \emph{Say it:} \emph{bhraṣṭācārada viruddha janaru bīdige iḷidaru.} — The people took to the streets against corruption.
  \item \emph{Recall:} say the Kannada for to promote, then the Kannada for to invest, then say \emph{bhraṣṭācāra} again
\end{itemize}

\begin{tcolorbox}[breakable,colback=teal!4,colframe=teal!35!black,title={Before you move on}]
What does \kn{ಭ್ರಷ್ಟಾಚಾರ} mean? (\textquotedblleft{}Corruption\textquotedblright{}.) Say it once more.
\end{tcolorbox}
\section[\texorpdfstring{\kn{ಲಂಚ} (lañca)}{ಲಂಚ (lañca)}]{\kn{ಲಂಚ} (lañca) — a bribe}

\label{lesson:KA-C402-lanca}



\begin{quote}
Before the new one: say the Kannada for to invest, then the Kannada for corruption.
\end{quote}

\subsection*{You'll want to know: \kn{ಲಂಚ}}
\textbf{\kn{ಲಂಚ}} — \emph{lañca} — \textquotedblleft{}a bribe\textquotedblright{}.

\begin{quote}
\kn{ಆ} \kn{ಅಧಿಕಾರಿ} \kn{ಲಂಚ} \kn{ತೆಗೆದುಕೊಂಡು} \kn{ಸಿಕ್ಕಿಬಿದ್ದನು}.
\end{quote}

\begin{quote}
\emph{ā adhikāri lañca tegedukoṇḍu sikkibiddanu.}
\end{quote}

\begin{quote}
That official was caught taking a bribe.
\end{quote}

\kn{ಲಂಚ} \kn{ಕೊಡು}, to bribe.

\begin{cousinweb}
\textbf{\kn{ಲಂಚ}} (\emph{lañca}) is borrowed from Sanskrit \emph{lañcā}, \textquotedblleft{}a present, a bribe\textquotedblright{} (CDIAL 10914).
\end{cousinweb}

\begin{grammarlens}[title={What you've built}]
One more word for giving reasons and opinions.
\end{grammarlens}

\subsection*{Your turn}
\begin{itemize}
\raggedright
  \item \emph{Say it:} \emph{lañca}
  \item \emph{Say it:} \emph{lañca}, once more
  \item \emph{Say it:} \emph{ā adhikāri lañca tegedukoṇḍu sikkibiddanu.} — That official was caught taking a bribe.
  \item \emph{Recall:} say the Kannada for to invest, then the Kannada for corruption, then say \emph{lañca} again
\end{itemize}

\begin{tcolorbox}[breakable,colback=teal!4,colframe=teal!35!black,title={Before you move on}]
What does \kn{ಲಂಚ} mean? (\textquotedblleft{}A bribe\textquotedblright{}.) Say it once more.
\end{tcolorbox}
\section[\texorpdfstring{\kn{ಪಾರದರ್ಶಕ} (pāradarśaka)}{ಪಾರದರ್ಶಕ (pāradarśaka)}]{\kn{ಪಾರದರ್ಶಕ} (pāradarśaka) — transparent}

\label{lesson:KA-C402-paradarsaka}



\begin{quote}
Before the new one: say the Kannada for corruption, then the Kannada for a bribe.
\end{quote}

\subsection*{You'll want to know: \kn{ಪಾರದರ್ಶಕ}}
\textbf{\kn{ಪಾರದರ್ಶಕ}} — \emph{pāradarśaka} — \textquotedblleft{}transparent\textquotedblright{}.

\begin{quote}
\kn{ಆಡಳಿತ} \kn{ಪಾರದರ್ಶಕವಾಗಿ} \kn{ಇರಬೇಕು}.
\end{quote}

\begin{quote}
\emph{āḍaḷita pāradarśakavāgi irabēku.}
\end{quote}

\begin{quote}
Administration should be transparent.
\end{quote}

Also of glass or water: \kn{ಪಾರದರ್ಶಕ} \kn{ಗಾಜು}.

\begin{cousinweb}
\textbf{\kn{ಪಾರದರ್ಶಕ}} (\emph{pāradarśaka}) is borrowed from Sanskrit \emph{pāradarśaka}, \textquotedblleft{}showing the opposite shore\textquotedblright{} (Monier-Williams): you see through to the far side.
\end{cousinweb}

\begin{grammarlens}[title={What you've built}]
One more word for giving reasons and opinions.
\end{grammarlens}

\subsection*{Your turn}
\begin{itemize}
\raggedright
  \item \emph{Say it:} \emph{pāradarśaka}
  \item \emph{Say it:} \emph{pāradarśaka}, once more
  \item \emph{Say it:} \emph{āḍaḷita pāradarśakavāgi irabēku.} — Administration should be transparent.
  \item \emph{Recall:} say the Kannada for corruption, then the Kannada for a bribe, then say \emph{pāradarśaka} again
\end{itemize}

\begin{tcolorbox}[breakable,colback=teal!4,colframe=teal!35!black,title={Before you move on}]
What does \kn{ಪಾರದರ್ಶಕ} mean? (\textquotedblleft{}Transparent\textquotedblright{}.) Say it once more.
\end{tcolorbox}
\section[\texorpdfstring{\kn{ಮರೆಮಾಚು} (maremācu)}{ಮರೆಮಾಚು (maremācu)}]{\kn{ಮರೆಮಾಚು} (maremācu) — to conceal}

\label{lesson:KA-C402-maremacu}



\begin{quote}
Before the new one: say the Kannada for a bribe, then the Kannada for transparent.
\end{quote}

\subsection*{You'll want to know: \kn{ಮರೆಮಾಚು}}
\textbf{\kn{ಮರೆಮಾಚು}} — \emph{maremācu} — \textquotedblleft{}to conceal\textquotedblright{}.

\begin{quote}
\kn{ಅವನು} \kn{ತನ್ನ} \kn{ತಪ್ಪನ್ನು} \kn{ಮರೆಮಾಚಿದನು}.
\end{quote}

\begin{quote}
\emph{avanu tanna tappannu maremācidanu.}
\end{quote}

\begin{quote}
He concealed his mistake.
\end{quote}

Past (I …): \kn{ಮರೆಮಾಚಿದೆ}, \emph{maremācide}.

What you hide takes -\kn{ಅನ್ನು}.

\begin{cousinweb}
\textbf{\kn{ಮರೆಮಾಚು}} (\emph{maremācu}) is native Kannada, a compound whose first part is \emph{maṟe}, \textquotedblleft{}concealment, a cover\textquotedblright{} (DEDR 4760); the second part, \emph{mācu}, is not settled.
\end{cousinweb}

\begin{grammarlens}[title={What you've built}]
One more word for giving reasons and opinions.
\end{grammarlens}

\subsection*{Your turn}
\begin{itemize}
\raggedright
  \item \emph{Say it:} \emph{maremācu}
  \item \emph{Say it:} \emph{maremācu}, once more
  \item \emph{Say it:} \emph{avanu tanna tappannu maremācidanu.} — He concealed his mistake.
  \item \emph{Recall:} say the Kannada for a bribe, then the Kannada for transparent, then say \emph{maremācu} again
\end{itemize}

\begin{tcolorbox}[breakable,colback=teal!4,colframe=teal!35!black,title={Before you move on}]
What does \kn{ಮರೆಮಾಚು} mean? (\textquotedblleft{}To conceal\textquotedblright{}.) Say it once more.
\end{tcolorbox}
\section[\texorpdfstring{\kn{ಶಿಫಾರಸು} \kn{ಮಾಡು} (śiphārasu māḍu)}{ಶಿಫಾರಸು ಮಾಡು (śiphārasu māḍu)}]{\kn{ಶಿಫಾರಸು} \kn{ಮಾಡು} (śiphārasu māḍu) — to recommend}

\label{lesson:KA-C402-sipharasu-madu}



\begin{quote}
Before the new one: say the Kannada for transparent, then the Kannada for to conceal.
\end{quote}

\subsection*{You'll want to know: \kn{ಶಿಫಾರಸು} \kn{ಮಾಡು}}
\textbf{\kn{ಶಿಫಾರಸು} \kn{ಮಾಡು}} — \emph{śiphārasu māḍu} — \textquotedblleft{}to recommend\textquotedblright{}.

\begin{quote}
\kn{ಮ್ಯಾನೇಜರ್} \kn{ಲಂಚ} \kn{ಪಡೆದು} \kn{ಅವನನ್ನು} \kn{ಕೆಲಸಕ್ಕೆ} \kn{ಶಿಫಾರಸು} \kn{ಮಾಡಿದನು}.
\end{quote}

\begin{quote}
\emph{myānējar lañca paḍedu avanannu kelasakke śiphārasu māḍidanu.}
\end{quote}

\begin{quote}
The manager took a bribe and recommended him for the job.
\end{quote}

Past (I …): \kn{ಶಿಫಾರಸು} \kn{ಮಾಡಿದೆ}, \emph{śiphārasu māḍide}.

For what: -\kn{ಗೆ}. \kn{ಶಿಫಾರಸು} alone is a recommendation.

\begin{cousinweb}
\textbf{\kn{ಶಿಫಾರಸು} \kn{ಮಾಡು}} (\emph{śiphārasu māḍu}) is built from \emph{śiphārasu}, a loan from Hindustani \emph{sifāriś}, \textquotedblleft{}a recommendation\textquotedblright{}, which comes from Persian, with \emph{māḍu}, \textquotedblleft{}to do\textquotedblright{}.
\end{cousinweb}

\begin{grammarlens}[title={What you've built}]
One more word for giving reasons and opinions.
\end{grammarlens}

\subsection*{Your turn}
\begin{itemize}
\raggedright
  \item \emph{Say it:} \emph{śiphārasu māḍu}
  \item \emph{Say it:} \emph{śiphārasu māḍu}, once more
  \item \emph{Say it:} \emph{myānējar lañca paḍedu avanannu kelasakke śiphārasu māḍidanu.} — The manager took a bribe and recommended him for the job.
  \item \emph{Recall:} say the Kannada for transparent, then the Kannada for to conceal, then say \emph{śiphārasu māḍu} again
\end{itemize}

\begin{tcolorbox}[breakable,colback=teal!4,colframe=teal!35!black,title={Before you move on}]
What does \kn{ಶಿಫಾರಸು} \kn{ಮಾಡು} mean? (\textquotedblleft{}To recommend\textquotedblright{}.) Say it once more.
\end{tcolorbox}
